\chapter{สูตรคณิตศาสตร์และเอกลักษณ์เวกเตอร์พื้นฐาน}

\section{ตัวดำเนินการในระบบพิกัดต่าง ๆ}
กำหนดให้ $\psi$ คือสนามสเกลาร์ และ $\vec{A}$ คือสนามเวกเตอร์

\subsection{\texorpdfstring{ระบบพิกัดคาร์ทีเซียน $(x, y, z)$}{ระบบพิกัดคาร์ทีเซียน (x, y, z)}}
\begin{align}
\nabla \psi &= \frac{\partial \psi}{\partial x}\hat{i} + \frac{\partial \psi}{\partial y}\hat{j} + \frac{\partial \psi}{\partial z}\hat{k} \\
\nabla \cdot \vec{A} &= \frac{\partial A_x}{\partial x} + \frac{\partial A_y}{\partial y} + \frac{\partial A_z}{\partial z} \\
\nabla^2 \psi &= \frac{\partial^2 \psi}{\partial x^2} + \frac{\partial^2 \psi}{\partial y^2} + \frac{\partial^2 \psi}{\partial z^2}
\end{align}

\subsection{\texorpdfstring{ระบบพิกัดทรงกระบอก $(\rho, \phi, z)$}{ระบบพิกัดทรงกระบอก (rho, phi, z)}}
\begin{align}
\nabla \psi &= \frac{\partial \psi}{\partial \rho}\hat{\rho} + \frac{1}{\rho}\frac{\partial \psi}{\partial \phi}\hat{\phi} + \frac{\partial \psi}{\partial z}\hat{k} \\
\nabla \cdot \vec{A} &= \frac{1}{\rho}\frac{\partial}{\partial \rho}(\rho A_\rho) + \frac{1}{\rho}\frac{\partial A_\phi}{\partial \phi} + \frac{\partial A_z}{\partial z} \\
\nabla^2 \psi &= \frac{1}{\rho}\frac{\partial}{\partial \rho}\left(\rho\frac{\partial \psi}{\partial \rho}\right) + \frac{1}{\rho^2}\frac{\partial^2 \psi}{\partial \phi^2} + \frac{\partial^2 \psi}{\partial z^2}
\end{align}

\subsection{\texorpdfstring{ระบบพิกัดทรงกลม $(r, \theta, \phi)$}{ระบบพิกัดทรงกลม (r, theta, phi)}}
\begin{align}
\nabla \psi &= \frac{\partial \psi}{\partial r}\hat{r} + \frac{1}{r}\frac{\partial \psi}{\partial \theta}\hat{\theta} + \frac{1}{r\sin\theta}\frac{\partial \psi}{\partial \phi}\hat{\phi} \\
\nabla \cdot \vec{A} &= \frac{1}{r^2}\frac{\partial}{\partial r}(r^2 A_r) + \frac{1}{r\sin\theta}\frac{\partial}{\partial \theta}(\sin\theta A_\theta) + \frac{1}{r\sin\theta}\frac{\partial A_\phi}{\partial \phi} \\
\nabla^2 \psi &= \frac{1}{r^2}\frac{\partial}{\partial r}\left(r^2\frac{\partial \psi}{\partial r}\right) + \frac{1}{r^2\sin\theta}\frac{\partial}{\partial \theta}\left(\sin\theta\frac{\partial \psi}{\partial \theta}\right) + \frac{1}{r^2\sin^2\theta}\frac{\partial^2 \psi}{\partial \phi^2}
\end{align}

\section{เอกลักษณ์ผลคูณเวกเตอร์}
\begin{align}
\vec{A} \cdot (\vec{B} \times \vec{C}) &= \vec{B} \cdot (\vec{C} \times \vec{A}) = \vec{C} \cdot (\vec{A} \times \vec{B}) \\
\vec{A} \times (\vec{B} \times \vec{C}) &= (\vec{A} \cdot \vec{C})\vec{B} - (\vec{A} \cdot \vec{B})\vec{C} \\
\nabla \times (\nabla \times \vec{A}) &= \nabla(\nabla \cdot \vec{A}) - \nabla^2 \vec{A}
\end{align}
